\subsection{Prediction Is Not Identification}
\label{sec:narrative-information}

For the exact observation map $\Psi$, compare a new intervention $e$ with
and without the same earlier history $h$:
\begin{equation}
 s_{e\mid h}(t)=\Psi_t(h\oplus e)-\Psi_t(h).
 \label{eq:history-increment}
\end{equation}
Both paths start from the same physical and instrumental state. For an additive intervention under the same linear time-invariant
dynamics, this increment is independent of the prior state. A topology
change can alter the state matrix, so that statement does not apply to
switched linear models in general. Neither switching nor nonlinearity
guarantees more source information.
Our stream implementation approximates both paths with local switched
predictors. It does not solve the full nonlinear DAE for every history.

A good fit to one source does not establish that competing sources can be
separated. With severity and nuisance coordinates $\beta_j$ profiled over
their admitted set $\mathcal A_j$, a local separation measure is
\begin{equation}
 d_h^2=\min_{j\ne h}\ \min_{\beta_j\in\mathcal A_j}
       \|s_h-s_j(\beta_j)\|_{\Sigma^{-1}}^2.
 \label{eq:narrative-separation}
\end{equation}
Small distance makes location noise-sensitive; it is neither a universal
error bound nor a calibrated posterior
\cite{eriksson2013diagnosability,scharf1994matched}. Fit, source ranking
and uncertainty calibration require separate tests. The supplementary
controls retain their known-onset and observation contracts; their state
reconstruction scores are not pooled with the causal study. Severity
fitting can absorb context error parallel to a source signature while
leaving intensity biased. This standard projection effect explains why
correct location and inaccurate response amplitude can coexist.
